MeBench makes experimentally supported regulatory dependencies an explicit evaluation target alongside expression outcomes. Among the released Task 2 records, 792 responses predict the correct direction but fail the mechanism chain (17.1\% of all responses). In the PPARA-activation subset, all 72 directions are correct while none of the chains passes. Under the released judgments, these cases expose a gap that endpoint accuracy leaves unresolved: whether the explanation preserves the dependency needed to interpret the intervention. This matters for biological hypothesis generation, where a proposed mediator and its function determine which follow-up experiment would test the explanation.

The module analyses localize this gap to distinct biological operations (Appendices~\ref{app:modules}--\ref{app:function}). The GATA1 cases distinguish preserved occupancy from impaired TAL1/LMO2-complex recruitment; the H2A.Z cases distinguish naming a mediator from assigning its context-specific repressive contribution. Cooperative-partner omissions add a third diagnostic target: recovering a required dependency when the queried intervention changes. Together, these contrasts show why mechanism evaluation should retain component identity, functional role and intervention logic as separate criteria. Their scope is the curated experiments and supplied inputs; they do not establish universal deficiencies across TF families or identify an architectural cause.

For model development, these diagnostics motivate supervision that connects mediator identity to its signed, context-specific role, together with paired intervention examples that require consistent use of the same dependency. Evidence-conditioned evaluation could then test whether supplying the relevant experiment repairs partner omissions, role inversions or both. Such studies should use source- and pair-disjoint tests, compare mediator-withheld and mediator-supplied conditions, and score endpoints and explanations separately. MeBench provides traceable cases for constructing these comparisons. Demonstrating a training benefit will require controlled model comparisons and independently adjudicated references; no such improvement is claimed here.
